\documentclass[10pt,a4paper,ragged2e,withhyper,normalphoto]{config/altacv}
\input{config/header_settings.tex}
\begin{document}

\name{Joe Egan}
  \tagline{High Energy Physics, Large Language Models, Data Science}

\personalinfo{ % these are all option, comment out any you don't want
  \email{joe.egan.23@ucl.ac.uk}
  \email{joseph.caimin.egan@cern.ch}
  % \phone{}
  \linkedin{joe-egan} % this just needs your <username>, which will look something like 'emily-smith', the link will render as https://www.linkedin.com/in/<username>
  % \homepage{www.ucl.ac.uk/astrophysics/name-phd-student}
  \github{joes-git} % this just needs your <username>, which will look something like 'emily-smith', the link will render as https://github.com/<username>
  \location{London}
}

\input{config/body_settings.tex}
\makecvheader
\begin{paracol}{2}

%% CDT Project Info
\cvevent{PhD: CDT Data Intensive Science, High Energy Physics}{University College London}{Oct 2023 -- Present}{London, UK}
\begin{itemize}
\item High Energy Physics student on the ATLAS experiment at CERN.
\item Physics research focuses on making model independent measurements with the ATLAS detector, and re-interpreting existing results to set novel constraints on physics beyond the Standard Model.
\item Experienced in building and deploying tools utlising Large Language Models, such as a Retrieval Augmented Generation (RAG) AI assistant for the ATLAS experiment.
\end{itemize}
%%

%% Experience
\cvsection[]{Experience}

\cvevent{chATLAS: AI assistant for the ATLAS collaboration}{The ATLAS experiment}{Oct 2024 -- Present}{CERN, Switzerland}
\begin{itemize}
\item Core developer on the chATLAS project, an AI assistant based on Retrieval Augmented Generation (RAG).
\item Built a modular suite of python packages to scrape ATLAS documentation, create and store document embedding vectors, efficiently search the database
 and pass the retrieved documents to a Large Language Model (LLM).
\item Built and deployed a scheduled CI/CD pipeline to periodically re-scrape the data sources and update the embedding database.
\item Paper in preparation.
\end{itemize}

\cvevent{Analysing Immigration Sentiment in UK politics with Machine Learning}{The Guardian}{Jan 2024 -- Apr 2024}{London, UK}
\begin{itemize}
\item Used natural language processing and large language models perform a quantitative analysis of immigration rhetoric in UK House of Commons speeches.
\item Built an end-to-end data analysis pipeline: speech scraping, increasing amount of relevant snippets for annotation, training a classifier, performing inference and analysis of results.
\item Engineered and tested prompts to use an LLM as an annotator. We were able to show that using an LLM did not degrade the inter-annotator agreement. This enabled the size of the training dataset for the classifier to be significantly increased with limited person power.
\item Paper in preparation.
\end{itemize}

\cvevent{Contur: Gaining new insights into physics Beyond the Standard Model from repurposed LHC results}{UCL}{Sep 2023 -- Present}{London, UK}
\begin{itemize}
\item Core developer on the Contur project, which reuses existing LHC measurements to constrain new physics models.
\item Overhauled the statistical calculations to use the Spey package, allowing for parameter fitting, modelling higher order signal processes and profiling nuisance parameters.
\item Applied Contur to set new limits on a Beyond the Standard Model theory involving dark matter candidates arising from the promotion of baryon number to a local gauge symmetry.
Paper published in Physical Review D.
\item Applied Contur to constrain the couplings of axion-like particles to top quarks and gluons. Submitted to JHEP.
\end{itemize}

\switchcolumn %back to the top, next column

%% Extracurricular Activities
\cvsection{Publications}

\begin{itemize}
  \item "Local baryon number at the LHC" \textit{Phys. Rev. D 112, 015012 (2025)} [DOI: 10.1103/2wbm-g6yv] [arXiv: 2505.06341]
  \item "Probing the coupling of axions to tops and gluons with LHC measurements"
  \item "The evolution of immigration political rhetoric in the  UK" \textit{In preparation}
  \item "chATLAS: An AI assistant for the ATLAS collaboration" \textit{In preparation}

\end{itemize}
%%

\cvsection{Selected Talks}

\cvevent
{ATLAS Week}
{chATLAS: An AI assistant for the ATLAS collaboration}
{Sep 2025}
{Paestum, Italy}

\cvevent
{ATLAS Software and Computing Week}
{Progress on chATLAS during my ATLAS Qualification Project}
{Oct 2025}
{Ljubljana, Slovenia}

\cvevent
{LHC Re-Interpretation Workshop}
{Constraints On New Theories Using Rivet: Version 3.1 update}
{Jan 2025}
{CERN, Switzerland}

%% Skills
\cvsection{Skills}
% \cvtag{A thing i am good at} 
% \cvtag{and another} \break
% \cvtag{ Lorem ipsum}
% \cvtag{rutrum quis mi.}

% \divider\smallskip

\cvtag{Python}
\cvtag{C++}
\cvtag{SQL}
\cvtag{Bash}
\cvtag{Git}
\cvtag{CI/CD}
\cvtag{uv}
\cvtag{ruff}
\cvtag{ty}
\cvtag{spaCy}
\cvtag{LangChain}
\cvtag{HPC}
%%

%% Language (this section is optional, if you only speak one language you can delete it , PS - GCSE is about A2/B1, A level is B2)
% \cvsection{Language}

% \cvskill{German \textcolor{body}{Native}}{5}
% \cvskill{English \textcolor{body}{C1}}{4}
% \cvskill{Language \textcolor{body}{level}}{2}
%%

%% Education
\cvsection{Education}

\cveduevent{MPhys (Hons.) Physics}{University of Manchester}{2018 -- 2022}{Manchester}{\nth{1} Class Hons (83\% average)}
\cveduevent{A Level Maths, Physics, Chemistry}{Silverdale School}{2011 -- 2018}{Sheffield}{A*A*A*}
%%

%% References
\cvsection{References}
\cvref{Prof. Jonathan Butterworth -- Supervisor}{j.butterworth@ucl.ac.uk}

\cvref{Dr. Gabriel Facini -- chATLAS Supervisor}{gabriel.facini@ucl.ac.uk}
\divider %do not leave a line between this and the last reference

Additional references available on request
%%

\end{paracol}
\end{document}
